\section{Imports y utilidades de Python}

\begin{lstlisting}[style=py]
import sys
from collections import Counter, defaultdict, deque
from functools import lru_cache
from heapq import heappop, heappush
from bisect import bisect_left, bisect_right
from math import gcd, isqrt, lcm, inf
# sys.setrecursionlimit(1_000_000) no se limite
input = sys.stdin.buffer.readline
\end{lstlisting}

\subsection{Entrada rapida}

\begin{lstlisting}[style=py]
n = int(input())
values = list(map(int, input().split()))
\end{lstlisting}

Lee una linea en $O(L)$, donde $L$ es su cantidad de bytes.

\subsection{Counter}

Cuenta frecuencias.

\begin{lstlisting}[style=py]
frequency = Counter([2, 1, 2, 3, 2])
print(frequency[2])       # 3
print(frequency.most_common(1)) # [(2, 3)]
\end{lstlisting}

Construirlo cuesta $O(n)$; consultar una frecuencia cuesta $O(1)$ promedio.

\subsection{defaultdict}

Crea un valor inicial automaticamente para claves nuevas.

\begin{lstlisting}[style=py]
graph = defaultdict(list)
graph[1].append(4)
graph[1].append(7)
print(graph[1]) # [4, 7]

count = defaultdict(int)
count[5] += 1
print(count[5]) # 1
\end{lstlisting}

Acceso e insercion: $O(1)$ promedio.

\subsection{deque}

Cola doble: insertar o quitar en ambos extremos.

\begin{lstlisting}[style=py]
q = deque([2, 3])
q.append(4)
q.appendleft(1)
print(q.popleft()) # 1
print(q.pop())     # 4
\end{lstlisting}

\texttt{append}, \texttt{appendleft}, \texttt{pop} y \texttt{popleft} cuestan
$O(1)$.

\subsection{lru\_cache}

Memoiza resultados de una funcion recursiva.

\begin{lstlisting}[style=py]
@lru_cache(None) # None almacena y si pongo nro ultimos x 
def fibonacci(n):
    if n <= 1:
        return n
    return fibonacci(n - 1) + fibonacci(n - 2)

print(fibonacci(10)) # 55
\end{lstlisting}

Cada estado se calcula una vez. En este ejemplo: $O(n)$ tiempo y $O(n)$ memoria.

\subsection{heapq}

Heap minimo: permite obtener siempre el menor elemento.

\begin{lstlisting}[style=py]
heap = []
heappush(heap, 5)
heappush(heap, 2)
heappush(heap, 7)

print(heappop(heap)) # 2
print(heappop(heap)) # 5
\end{lstlisting}

\texttt{heappush} y \texttt{heappop} cuestan $O(\log n)$.
Para un max-heap, guardar los valores negados.

\subsection{bisect}

Busca posiciones en un vector ordenado.

\begin{lstlisting}[style=py]
values = [1, 2, 2, 2, 5]

print(bisect_left(values, 2))  # 1
print(bisect_right(values, 2)) # 4
print(bisect_left(values, 4))  # 4
\end{lstlisting}

\texttt{bisect\_left} busca la primera posicion posible; \texttt{bisect\_right},
la posterior a las repeticiones. Ambos cuestan $O(\log n)$.
Insertar en una lista sigue costando $O(n)$.

\subsection{Matematica}

\begin{lstlisting}[style=py]
print(gcd(18, 24)) # 6
print(lcm(6, 8))   # 24
print(isqrt(20))   # 4
print(inf > 10**100) # True
a = int("ABCDEF", 16)
b = int("123456", 16)

result = a * b
print(format(result, "X"))
\end{lstlisting}

\texttt{gcd} cuesta $O(\log(\min(a,b)))$; \texttt{lcm} tiene costo equivalente.
\texttt{isqrt(n)} devuelve $\lfloor\sqrt n\rfloor$ sin errores de punto flotante.
\texttt{inf} representa infinito para inicializar minimos o distancias.